% =============================================================================
%  Upper Volume: Articles 3--15
%  Generated from translation/pages_19-22.md, pages_23-26.md, pages_27-30.md
% =============================================================================

% --- Article 3 is in a prior translation file (pages_15-18.md) ---
% --- This fragment begins at Article 4 ---

\article{4}{On the Resolve to Plunge into the Sea When Everyone Shares the Same Outlook}

\begin{sokyubox}
{\small\itshape\color{sokyubar} \jp{第四章\quad 人も同じ見込の節海中へ飛入る心持の事}}

\medskip
When rice has been falling steadily, when there is no change in the
\jt{上方}{Kamigata} market, when rumours abound that domain rice
(\jt{拂物}{haraimono}) from the provinces and from \jt{最上}{Mogami}
is plentiful, when market sentiment (\jt{人氣}{ninki}) is uniformly
bearish, and when even one's own judgment inclines toward weakness with
no telling how far prices may fall---at precisely that moment, reverse
one's thinking and buy. This resolve requires the spirit of plunging
into the sea. It is an exceedingly unpleasant feeling, yet at that
moment one must not give rise to doubt: buy. This is supremely
profitable. When one expects a decline and the market obliges, it is
easy enough; but when sentiment leans entirely toward decline, the
market perversely rises, and one's calculations avail nothing. The same
applies in the opposite direction. This spirit of plunging into the sea
is the innermost secret.
\end{sokyubox}

\commentary
Mogami refers to the collective name for what is now the Murayama and
Mogami districts. \textit{Haraimono} denotes the official disposal rice
(\jt{御拂米}{ohara-mai}) of the various domains, all of which was
customarily shipped to the port of \jt{酒田}{Sakata} for sale---a point
already noted at the opening of this volume.

The previous chapter explained selling when sentiment has swung
unanimously to the buy side; this chapter explains buying when sentiment
has swung unanimously to the sell side.

When sentiment has tilted entirely toward selling and one's own
sympathies likewise favour the sell side, one must absolutely not sell.
At such a time, one should summon the great courage of one who plunges
into the sea and turn to the buy side. The reason is this: the
phenomenon of unanimously bearish sentiment, seen from one angle, means
that sellers have sold to the absolute limit---the season when
profit-taking (\jt{利喰}{rigui}) must soon emerge is at hand. And since
everyone, without exception, has piled onto the sell side, the subtle
mechanism of a reactive upward turn is already in motion. However, what
the master teaches here applies to cases where both the number of months
and the price level have ripened. It must not be applied when the months
and days are still young. \jt{慈雲齋}{Jiunsai} also says as much:

\textit{``Whether rising or falling, ride the move at five parts or
one-tenth; know that at two- or three-tenths, the market reverses.''}

\textit{``Whether rising or falling, at five parts or one-tenth the ride
is good; past the midpoint, riding is foolish.''}

Together with this chapter, these are worthy admonitions.

% --- Article 5 ---
\article{5}{On Rice That Consolidates from Winter through the First and Second Months}

\begin{sokyubox}
{\small\itshape\color{sokyubar} \jp{第五章\quad 冬より正二月迄持合ふ米の事}}

\medskip
Rice that holds in consolidation (\jt{持合}{mochai}) at the floor price
from mid-winter through the first and second months will, from the third
and fourth months into the fifth and sixth, assuredly rise.
\end{sokyubox}

\commentary
When prices were high the previous year and then declined for three
months or more through January of the following year, this is a market
that fell because the previous year's harvest was abundant, the autumn
supply was ample, and bearish drift (\jt{ダラダラ}{daradara}) was
chanted to exhaustion. With no immediately bullish factors in sight, and
particularly when the winter weather is unobjectionable, sellers appear
early in anticipation of fine equinox weather, and for a time low prices
will persist. This low price often becomes the year's floor price
(\jt{底直段}{soko nedan}). One should therefore consult the chart
(\textit{keisen-hy\=o}), examine the number of months elapsed and the
proportional extent of the decline, weigh these against external
conditions in the broader economy, and resolve to buy. This is what is
called the buying opportunity (\jt{買時機}{kai-jiki}). The reason is
that in such calm, uneventful years, the market has already exhausted
whatever decline it could produce. Even at the equinox---a time when
prices ought to fall---they do not fall; and whenever the weather turns
the slightest bit unfavourable, buying driven by hasty profit-taking and
the like emerges, gradually lifting prices. On top of this, as the sixth
month approaches and the planting season presses near, if the weather is
unpromising, short-covering by sellers and fresh buying appear, producing
unexpectedly high prices. This sometimes marks a temporary ceiling
(\jt{一ト天井}{hito-tenj\=o}), so one must study the number of months
and the proportional move to determine whether it is indeed the ceiling
price.

% --- Article 6 ---
\article{6}{On Markets That Fall Sharply and Rise Sharply}

\begin{sokyubox}
{\small\itshape\color{sokyubar} \jp{第六章\quad 急に下げ急に上る相塲の事}}

\medskip
In a market that falls sharply and rises sharply, the time limit for the
ceiling and floor cannot be fixed in advance. One must judge the moment
and close out the trade (\jt{手仕廻}{teshimai}). At such times: ``At
two, close out; at three, fully ripe; at four, it turns''---this is the
secret tradition of the Three Positions (\jt{三位の秘傳}{sanmi no
hiden}).
\end{sokyubox}

\commentary
The best standard for determining the ceiling and the floor is to
observe whether the number of months and the proportional move have
ripened or not. However, there is an exceptional case: namely, the
market described in this chapter---one that rises sharply or falls
sharply. This is what is called a fast-ripening market
(\jt{早熟相塲}{s\=ojuku s\=oba}). It reaches its ceiling or its floor
without regard to the number of months or days. Therefore, if a trade
(\jt{商内}{akinai}) shows a profit, one should close it out at a
reasonable point. The guiding principle is to keep in mind the rhythm of
``at two, close out; at three, fully ripe; at four, it turns''---release
one's position, observe with a calm and level mind as the market ripens
fully, and then strike at the reversal.

This chapter explains the approach to markets that form an exception to
the calendar rule.

% --- Article 7 ---
\article{7}{On How to Read Rice at the Ceiling Price from Winter through the New Year}

\begin{sokyubox}
{\small\itshape\color{sokyubar} \jp{第七章\quad 冬より正月迄天井直段の米見様の事}}

\medskip
Rice at the floor price through mid-winter and the New Year should rise
by the fifth and sixth months.

Rice at the ceiling price from mid-winter through the first and second
months should fall by the fifth and sixth months. When it falls fully in
the fifth month, it will surge upward in the sixth. If it does not fall
in the fifth month, it will assuredly break in the sixth---of this there
is no doubt.

Rice at the floor price through the seventh, eighth, ninth, and tenth
months should be expected to rise by the twelfth month.
\end{sokyubox}

\commentary
This chapter should be read in three sections. The first section has
already been explained in detail in the preceding chapter, so we proceed
directly to the second section.

Rice that has continued rising from the previous year owing to poor
harvest and reached its ceiling price by around the twelfth or first
month: even after physical rice from the supplying regions, drawn by the
high price of the futures contract, has come to market in fair volume,
sentiment remains buoyant because people remember the previous year's
crop failure and continue buying on nostalgia. It does not collapse
easily. Nevertheless, in such circumstances one must not forget the sell
policy appropriate to a ceiling-price consolidation (\jt{天井保合}{tenj\=o
mochai}) following a poor harvest. The reason is that such nostalgic
buying cannot last forever. By the fifth month, the wheat harvest is
average, the weather is not unfavourable, and moreover the poor-quality
rice from the previous year's bad crop risks spoilage in the warm
weather. The rush of rice-holders to sell, among other factors, produces
a major break.

However, when the decline is excessive, a sharp upward reaction may
occur into the sixth month, so one should take note. Next, let us
explain the third section.

A floor price in the eighth or ninth month arises in a year when the
weather since spring has been unobjectionable, when the wheat harvest has
been proclaimed bountiful, and when prices have gradually declined.
Occasionally a rumour of flooding in one or two provinces produces a
small rebound, but on the whole, when post-planting weather has been good
and the midsummer heat (\jt{土用}{doy\=o}) has been ample, farmers sell
their stored rice with confidence and merchants adopt a policy of
refraining from buying, making it difficult for purchases to advance.
Prices drift lower into the \jt{盆}{Bon} season; and as everyone,
anticipating an abundant harvest (\jt{豐年萬作}{h\=onen mansaku}), joins
in chanting bearish drift, the year's floor price sometimes appears. At
such a juncture, one should decidedly buy. Because sentiment is leaning
uniformly toward the bearish side, a minor storm warning of low intensity
will scarcely disturb confidence. But while things stand thus, a rumour
of insect damage in a nearby province emerges, or a typhoon strikes, or
an early cold spell arrives---some trivial event becomes the catalyst,
and prices gradually advance. By around the eleventh month, word spreads
of insufficient sickle-entry (\jt{鎌入不足}{kamairi-fusoku}---meaning
inadequate harvesting), or of unsettled autumn-drying weather, and prices
climb ever higher. On top of that, complaints that producing regions find
the futures price unprofitable, or that producers are withholding stock,
and various other circumstances converge, tightening the market
relentlessly. Short-covering by sellers (\jt{煎米}{iri-mai}) and bullish
buying pressure drive prices naturally upward through the eleventh and
twelfth months. However, a market that has risen on these circumstances
often reaches a temporary ceiling around the twelfth or first month, so
great attention is required.

% --- Article 8 ---
\article{8}{On Ceilings in the Seventh through Tenth Months, and the Summer Decline}

\begin{sokyubox}
{\small\itshape\color{sokyubar} \jp{第八章\quad 七、八、九、十月天井及夏下げの事}}

\medskip
Rice at the ceiling price in the seventh, eighth, ninth, or tenth month
should be expected to decline by the twelfth month. In a year when
speculative positions (\jt{思入}{omoiire}) are large, the summer decline
(\jt{夏下げ}{natsu-sage}) will assuredly occur.
\end{sokyubox}

\commentary
Generally speaking, a ceiling price in the seventh, eighth, ninth, or
tenth month is, in most cases, the result of natural disasters such as
storms and floods during the summer, or else a temporary high produced by
the thinning of supply at the crop boundary (\jt{端境}{hazakai}).
Thereafter, once the weather recovers and crop prospects are reassessed
upward, and as the high prices previously attracted physical rice from
various regions, futures too gradually decline. Then, as the harvest
approaches and the prospect of ample new-crop intake becomes clear,
farmers in the countryside begin to contract forward sales of new rice at
low prices to merchants. Physical rice in the various regions falls
first, and the resulting impact on futures---including liquidation by the
bull side---produces an unexpected decline. The subsequent phrase, ``In a
year when speculative positions are large, the summer decline will
assuredly occur,'' carries the same meaning as the second section of the
preceding chapter.

\begin{sokyubox}
{\small\itshape\color{sokyubar} Article 8 (continued)}

\medskip
This is the fourth sub-section of the chapter, continuing from the
previous one. It explains the pattern of the rise from the floor price
(\jt{底直段}{soko nedan}) to the ceiling, and then describes the
consolidation (\jt{保合}{mochai}) at the ceiling in the spring following
a poor-harvest year.
\end{sokyubox}

To illustrate the pattern of a market rising from the year's floor
price: this market falls five tawara (\jt{俵}{tawara}) and then rises
ten, falls ten tawara and then rises twenty---oscillating upward in this
fashion, pressing higher from around the eighth month through the first
month of the following year, grinding upward for five or six months until
it produces the ceiling price. In an ordinary year, when this ceiling is
struck, the market drops sharply. In an extreme poor-harvest year,
however, sentiment is so taut that the market may hold in consolidation
near the ceiling for two or three months without falling. Even so, from
around the ceiling month (the first month) and the month following (the
second month), the tone softens slightly, and by the fourth or fifth
month---because rice is dear and therefore unprofitable---buyers
(referred to as ``ships'': Sakata being a port town, buyers were largely
from the shipping trade) grow reluctant. Then, when the sixth month
brings fair weather and the midsummer heat (\jt{土用}{doy\=o}) bears
down, sentiment turns worse. In particular, the sixth month brings anxiety over
rice spoilage and measure-loss (\jt{桝減り}{mashiberi}), and it is also
the opening month (\jt{發會}{hakkai}) of the new November contract
(\jp{新甫十一期}); for all these reasons, amateurs (\textit{shir\=oto})
and out-of-town long holders alike rush to close their positions, and the
market crashes seventy or eighty to as much as a hundred tawara.

% --- Article 9 ---
\article{9}{On Maneuvering in a Poor-Harvest Year}

\begin{sokyubox}
{\small\itshape\color{sokyubar} \jp{第九章\quad 不作年駈引の事}}

\medskip
When in this region the sixth and seventh months are rainy and cool, the
weather persistently chilly, and clear skies rare, the harvest in this
area and neighbouring provinces will be extremely poor. Moreover, in
Kyushu, the western provinces, the \jp{中國} region, the Home Provinces
(\jt{五畿内}{Gokinai}), the \jp{東海道}, and the northern interior
(\jt{奥筋}{Okusuji}), weather and crop conditions differ from year to
year. The northern provinces may have a good harvest while the western
provinces fail, or the western provinces may prosper while the Kanto
fails. Although in most years conditions roughly follow those of one's
own province, there are exceptions. Consider this carefully.
\end{sokyubox}

\commentary
This chapter can be divided into several sections. The present section
describes the method of using local crop conditions as a basis for
surveying the national harvest outlook. Because the explanation here
takes the \jt{荘内}{Sh\=onai} region as its base point, ``this region''
(\jt{當地}{t\=ochi}) should be understood accordingly.

In general, the sixth and seventh months of the lunar calendar are an
extremely critical season for the growth of rice plants (\jt{稻草}{inakusa}).
Abundant sunshine and scorching heat are required. If rainfall is
frequent and the air feels cool---the kind of unseasonable weather where
one wears padded garments during the midsummer \textit{doy\=o}
period---then one should know that the harvest in this area and the
neighbouring provinces will be extremely poor. Furthermore, although it
is the same Japan, differences in the orientation of mountains and seas,
in ocean currents, and in meteorological conditions mean that weather
patterns and consequently crop yields vary to some degree among Kyushu,
the western provinces, Chugoku, the Home Provinces, the Tokaido, and the
northern provinces---they are not uniform. And so, when crop reports come
in from various regions, conflicting and contradictory, one sometimes
finds it difficult to judge which to follow. The essential point is to
use one's own province's crop conditions as the benchmark, fix the
general picture in one's own mind, and remember that---northern provinces
good harvest, western provinces poor; or western provinces good, Kanto
poor---there are always year-to-year differences.

\begin{sokyubox}
{\small\itshape\color{sokyubar} Article 9 (continued)}

\medskip
Furthermore, even when this region has had unsettled weather in the sixth
and seventh months, and the rice stalks are short, the paddy surfaces
appear concave, and the base growth looks thin, if the sun shines
continuously from the end of the sixth month through around the twentieth
of the seventh month, the outlook rapidly improves and shifts toward a
good harvest. Moreover, between the sixth and eighth months, pay the
closest attention to the severity of natural disasters---typhoons,
floods, and insect infestations. This applies not only to this region
but above all to Kyushu and the Kamigata region.
\end{sokyubox}

This section is the second part, explaining how to read the condition of
the rice plants. Now then: when unsettled weather persists in the sixth
and seventh months, and looking out over the paddies one sees that the
rice stalks are short, that the fields appear concave, and that the base
growth looks thin (in a bountiful year the paddy surface appears convex),
if sunshine continues through the end of the seventh month, the outlook
rapidly improves toward a good harvest. Also, the period from the sixth
through the eighth month is called the natural disaster season, a
critical time for the growth of the rice plants. Not only in this
region, but in all provinces, one must pay careful attention to the
occurrence and extent of natural disasters.

\begin{sokyubox}
{\small\itshape\color{sokyubar} Article 9 (continued)}

\medskip
Now, the November contract (\jt{霜月限}{shimotsuki-giri}) new trading:
taking into account old rice and reports of the year's crop, trading
opens (\jt{發會}{hakkai}) fifty to sixty or up to one hundred tawara
(\jt{俵}{tawara}) below the old-rice price. From that opening, a
further decline of forty or fifty tawara is rare. Because trading opens
based on the year's estimated harvest, in most years the price settles
just ten to twenty or thirty tawara lower, and depending on the year's
developments it turns upward. If the Kamigata region or this area
suffers poor harvests or natural disasters, a sharp rise may occur from
the time trading opens in the sixth month. In any event, each year the
severity of the crop failure, the quantity of stored old rice, and
conditions in Kyushu and the various provinces determine the price
movements in \jt{大坂}{\=Osaka}, which then produce corresponding
movements here. One must not be complacent. Truly, the variations are
without limit.
\end{sokyubox}

\commentary
This is the third section, describing the normal pattern of new trading
each year. The November contract new trading (\textit{shimotsuki-giri
shin-akinai}) refers to the November-delivery contract. By long custom
at the \jt{酒田米會所}{Sakata Kome Kaisho}, trading in the new-crop
November contract was opened from the first or last ten days of the
sixth month.

Now, as for the pattern of the market at its formation: in an ordinary
year, based on the wheat harvest outlook and the state of the growing
rice plants, the November contract opens at a discount of roughly fifty
to sixty tawara per hundred \textit{ry\=o}---that is, five to six or up
to ten tawara per ten-\textit{ry\=o} unit---below the old-rice price.
However, it is extremely rare for the price to fall a further four or
five tawara from its opening level. In most cases, after declining one,
two, or three tawara, it finds a floor and turns upward. The reason is
that the sixth month of the lunar calendar is not only the onset of the
natural disaster season, but also because the rice plants' growth
outlook, to the extent it is good, has already been discounted into the
low opening price---hence the price cannot be driven much further down.

\begin{sokyubox}
{\small\itshape\color{sokyubar} Article 9 (continued)}

\medskip
When the ceiling price emerges in mid-winter, it is because both the
local region and the \jt{上方}{Kamigata}---along with neighboring
provinces---have suffered natural disasters or crop failure from cold
weather. Both Osaka and the local market have risen to the point where
no one knows how much further they may climb. After five or six months
of grinding higher, with ten, twelve, or thirteen consecutive days of
people clamoring and the atmosphere crackling with tension---this is the
year's full ceiling price. One should regard it with alarm and never let
one's guard down. Sell with the resolve of one leaping into fire; great
profit is beyond doubt. Though this is exceedingly difficult to execute,
if one rolls one's position forward month after month for five or six
months, great profit is beyond doubt. After this ceiling has appeared, it
is well to observe the market for about ten days until it settles, and
then sell into it. Apart from this, when one establishes a speculative
sell position at other times, unexpected rallies invariably come and
losses are certain. One should always be on guard against this.
\end{sokyubox}

\commentary
This sub-section explains the proper mindset for selling at the ceiling
price. ``Natural crop failure'' (\jp{自然の不作}) refers to poor
harvests caused by cold weather.

To describe what it looks like when the ceiling price emerges in
mid-winter: beforehand, both the local region and the Kamigata have
experienced typhoons and other natural disasters, or cold weather has
produced crop failure. As a result, both Osaka and the local market have
been rising steadily, grinding upward for five or six months, until
sentiment reaches a state where no one knows how much further prices may
climb. On top of this, ten or twelve or thirteen more days of advances
pile on, and bulls and bears alike join together in a frenzy. When this
happens, it is the year's absolute highest price---the great ceiling
(\jt{大天井}{\=otenj\=o})---and the finest possible selling
opportunity. One should regard it with alarm and never let one's guard
down. The instruction is to sell with the resolve of one leaping into a
fire. After the great ceiling has been struck, the market will continue
falling month after month for five or six months; therefore, if one
rolls one's position at intervals over that span, great profit is beyond
doubt. That said, to sell into the great ceiling itself requires either a
thoroughly seasoned trader or extraordinary resolve, and so the text
advises: after the ceiling, observe carefully for five or six days, up to
ten days, and then sell into the second peak---the so-called ``second
mountain'' (\jt{二の山}{ni no yama}) in common parlance.

In truth, the only sure-win opportunity is the case where the market has
overshot---the ceiling price has run to a point out of line with all
other regions. Markets are, by their nature, eight parts bullish, and
therefore establishing speculative sell positions during a consolidation
(\jt{持合相場}{mochai s\=oba}) is inadvisable: unexpected rallies come,
and losses are frequent. One should keep this in mind at all times.

\begin{sokyubox}
{\small\itshape\color{sokyubar} Article 9 (continued)}

\medskip
From the first appearance of new rice, the market rises through
range-bound oscillation (\jt{通高下}{kayoi takasage}); therefore, when
rice is trending upward, one must not reckon with the abacus. Buy on
each dip within the oscillation, target a hundred-tawara rise, and
understand that a margin call (\jt{落引}{ochibiki}) of about two rounds
marks the limit of one's buying.
\end{sokyubox}

\commentary
This sub-section teaches that once the floor price of new rice has been
identified, one should adopt a buying stance and hold to it through the
intervening oscillations without wavering, and moreover, that when the
market has taken on an upward momentum, one must never trade against it.

In essence, the new contract (\textit{shinsho}) rises gradually from
around the opening (\jt{發會}{hakkai}) through range-bound oscillation.
This is termed the ``natural rise.'' Because such a market leads physical
rice on the strength of prevailing conditions, one must not compare
prices against other regions, nor measure the spread against physical
rice and complain of unprofitability---one must not sit with the abacus
and calculate. It is reasonable to expect a rise of roughly twice the
margin deposit (\jt{證據金}{sh\=oko-kin}); buy on each dip and
accumulate. In the course of such a rise, breakaway gaps will also
narrow, so one may continue accumulating even if margin calls come twice
on the losing side.

The phrase ``about two rounds of margin calls'' means being called for
additional margin (\jt{追證據金}{tsui-sh\=oko-kin}) twice. For the full
details, one should consult the section on margin calls in the ``Customs
of Sakata Rice Trading'' at the front of the volume.

\begin{sokyubox}
{\small\itshape\color{sokyubar} Article 9 (continued)}

\medskip
However, in a year when the harvest across the provinces is reasonable
and the local crop is about eight parts in ten, the ceiling may come
after a rise of only sixty or seventy tawara, depending on the year. As
a rule, a hundred-tawara rise is the norm. But when Kyushu, the
Kamigata, and the local region together suffer a harvest of four to five
parts---a half-crop year---market sentiment heats up exceptionally.
Visiting traders (\jt{旅人}{tabibito}) are a given, and everyone piles
into speculative buying. Moreover, in such a poor-harvest year, farmhouse
rice (\jt{内米}{uchi-mai}) is insufficient; even rice normally earmarked
for domain tribute payments (\jt{郷方札納}{g\=okata satsu-n\=o}),
provisions (\jt{夫食}{bujiki}), and other normally uncommitted stocks
gets bought up. The market is bid up without any reckoning until there is
no arbitrage rice (\jt{乘米}{nosei-mai}) left relative to Osaka. One
should pay close attention. In such cases, the rise is not limited to a
hundred tawara but may reach a hundred and twenty or thirty, even a
hundred and eighty tawara. Judging the moment is critical. Once the
ceiling has appeared, one must not hold a long position no matter how
convinced one may be. Should one mistakenly insist on being bullish, the
loss will assuredly be ruinous---enough to overturn one's entire fortune.
This is the state where sentiment has climbed to the absolute extreme of
bullishness: a hundred out of a hundred people believe prices will rise
without limit, the fever pitch of euphoria carries the crowd, and the
bulls multiply still further. One must exercise the strictest caution.
\end{sokyubox}

\commentary
``A half-crop year of four or five parts'' means a year when the harvest
yields only four parts or about half the normal crop. ``Visiting
traders'' (\jp{旅人}) means traders who have come from various
provinces---that is, out-of-town buyers. ``Farmhouse rice''
(\jt{内米}{uchi-mai}) is rice held at the farm itself. ``Domain tribute
payments'' (\jp{郷方札納}) is the tax rice farmers deliver to the domain
government. All of these terms are explained in the ``Customs of Sakata
Rice Trading'' at the front of the volume.

That said, when the harvest across all provinces is roughly average and
the local crop is about eight parts, the ceiling may come after a rise of
only sixty or seventy tawara; but as a rule, a hundred-tawara rise is the
standard. If, however, the Kamigata and the local region together produce
something close to a half crop, then buying from visiting traders
intensifies, and speculative players emerge in the local market as well.
In such a poor-harvest year, farmers' harvested rice is meager, so some
even purchase physical rice from outside sources to fulfill their tribute
obligations to the domain lord; moreover, purchases for famine-reserve
stockpiles (\jt{備荒貯蓄米}{biko chochiku-mai}) and other buying that
is entirely unnecessary in normal years also appear. These bullish
factors accumulate one upon another, stimulating market sentiment
until---measured against the Osaka price---the spread vanishes entirely or
even inverts, with prices bid up beyond all calculation or reckoning. One
must never neglect caution. And when the ceiling has appeared in such a
market, one must absolutely not hold a long position, no matter how
promising the outlook may seem. Should one insist on being bullish and
try to hold on, one may encounter an unexpected crash and suffer losses so
great as to endanger one's entire estate. Why must one not buy this rice?
Because in these cases, sentiment has climbed to the absolute extreme of
bullishness---a hundred people out of a hundred believe that prices will
rise without limit---and the bulls proliferate further still, and people
are easily beguiled by the frenzy into holding on. But such a moment is
the very apex of sentiment, and it is often the great ceiling for the
entire year.

\begin{sokyubox}
{\small\itshape\color{sokyubar} Article 9 (continued)}

\medskip
When rice that has risen from the floor is bought and a profit of twenty
or thirty tawara per unit accrues, there are those who sell it back for a
profit-taking (\jt{利喰}{rigui}). This is a grave error. One must
absolutely not sell. If one does sell back, then by the time the November
contract (\jt{霜月}{shimotsuki}) expires and the ceiling price appears,
cascading margin calls (\jt{連落}{tsure-ochi}) will result, leaving one
with no gain---or worse, producing actual losses. This is a critical
point. When the market approaches the ceiling, the first priority is to
ascertain the margin calls. For positions where the counterparty has been
called out, one may sell back. The above pertains to a severe
poor-harvest year of four to five parts.
\end{sokyubox}

\commentary
This sub-section teaches that once one has identified the floor and
adopted a buying stance, one should not waver at each range-bound dip,
but rather let one's profits run to their fullest extent. This is the
same counsel as David Ricardo's dictum: \textit{``Let your profits
run.''}

To explain: when rice is bought near the floor and a mere twenty or
thirty tawara of profit per contract accrues, some may engage in
profit-taking. But this is a serious error. The reason is that trading
custom at the \jt{酒田米會所}{Sakata Kome Kaisho} at the time did not
settle profit-and-loss on a position until the contract month expired,
except for positions where the counterparty had been called on margin.
Therefore, even if one sold back a long position at a gain of two or
three tawara, should the market subsequently rise further, the sold-back
portion would be subject to margin calls, and the resulting cascading
calls would not merely eliminate one's gains but actually produce losses.
(See ``Customs of Sakata Rice Trading.'') Therefore,
when one judges the market to be near the ceiling, the essential thing is
to examine the state of one's counterparty's positions. For positions
where margin calls have been triggered, the profit is realized
automatically in the settlement, so there is nothing further to consider.
But if the counterparty has posted additional margin and is carrying
(\jt{繋ぎ}{tsunagi}) the position, then one must sell against that
portion as a hedge. However, the maneuvering (\jt{駈引}{kakehiki})
described above pertains to the tactics of a four-parts or half-crop
poor-harvest year.

% --- Article 10 ---
\article{10}{On the Great Frenzy of a Rising Market}

\begin{sokyubox}
{\small\itshape\color{sokyubar} \jp{第十章\quad 相塲引上げ大騒考の事}}

\medskip
From the very start of the new contract, the market steadily advances, a
great frenzy erupts, and the ceiling price is struck. Suppose on that
day the price reaches two tawara three bu, hits the ceiling, and within
the same day reverses wildly to three tawara two or three bu. Pay close
attention to that day's two tawara three bu. Within a day or two, the
price will oscillate back to two tawara five or six bu, but it will lack
the momentum to reach two tawara three bu and will drop limply
(\textit{kurari}) back to the three-tawara range. This is where one
should sell. Because the rice was rapidly pulled up to two tawara three
bu, sentiment remains taut within the three- to four-tawara range; when
it falls to the four-tawara range, buyers push it back to the
three-tawara level. But gradually buyers thin out, and the price drops
to five or six tawara. At this point, one should consider arbitrage rice
and buy in---the market will return to the three-tawara range at most.
This is the ``great oscillation'' (\jt{大通ひ}{\=ogayoi}). If, in the
following month, the price surpasses two tawara three bu by even one or
two bu, buy without hesitation. If, in the month after that, it reaches
only about three tawara to two tawara six or seven bu and stalls there,
sell aggressively on conviction. This is the secret of the
Three-Position Tradition (\jt{三位の秘傳}{sanmi no hiden}).
\end{sokyubox}

\begin{editorialnote}
The price figures here use the Sakata exchange's bale-and-bu notation.
``Two tawara three bu'' is the highest price, with lower numbers
representing higher prices on an inverted scale relative to modern
conventions. Hayasaka's commentary below restates the principle using
hypothetical yen-and-sen figures for clarity.
\end{editorialnote}

\commentary
To illustrate the great oscillation following the first ceiling: after
the opening of the November contract, the market grinds steadily upward
until---let us suppose---in the morning session it posts a high of
ninety-five sen, and within that same day falls back to seventy-five sen.
Pay very close attention to that day's ninety-five sen. Within a day or
two, the price will oscillate from eighty-five or eighty-six sen up to
ninety or ninety-one or ninety-two sen, but it will lack the momentum to
reach ninety-five, and it will drop limply back to the seventy-sen range.
This is where one should sell. Because the rice surged to a high of
ninety-five sen and showed ninety-one or ninety-two, even when the price
falls through the seventy-sen range to the sixty- and fifty-sen range,
sentiment remembers the high price, and buy-backs are plentiful; when it
drops to the fifty-sen range, it is pushed back up to seventy or eighty
sen. But gradually the buying thins out, and the price falls to thirty
or forty sen, even twenty sen. At this point, one should weigh the
spread against other regions, the comparison with physical rice, and take
profit. The market will assuredly rally again to the seventy- or
eighty-sen range---this is what is called the ``great oscillation.''
However, this is merely the great oscillation \textit{after the first
ceiling}, not the true ceiling. If, in the following month, the price
surpasses that ninety-five-sen monthly high (\jt{月節}{tsukibushi})---the
previous peak---and pushes into the full-yen range, one should buy
without hesitation. And if, in the month after that, the price merely
reaches seventy or eighty sen and stalls, one should sell aggressively.
This, the author says, belongs to his Three-Position Tradition. (The
prices in the above explanation are hypothetical figures used for
convenience of illustration; one should not take them literally.)

% --- Article 11 ---
\article{11}{On Prices That Arrive at the Ceiling}

\begin{sokyubox}
{\small\itshape\color{sokyubar} \jp{第十一章\quad 行付天井直段の事}}

\medskip
What is meant by ``prices arriving at the ceiling''
(\jt{行付天井直段}{yukitsuki tenj\=o nedan}) is this: over the course
of several months, the market swings up and down, gradually pulling
prices up by fifty or sixty tawara. Thereafter a period of uncertainty
follows---it is unclear whether prices will rise further or fall. Then,
day after day, prices are pulled up relentlessly by ten tawara at a
time; both buyers and sellers are swept along; and when great commotion
erupts, prices peak. This is the annual ceiling price.
\end{sokyubox}

\commentary
This chapter describes the pattern of market sentiment when the ceiling
price is struck.

What is called the ``ceiling-arrival price'' (\jp{天井行付直段})---also
termed \textit{yukitsuki tenj\=o}---refers to the highest ceiling of the
entire year. Now, the pattern of sentiment when this ceiling forms is as
follows: over a span of several months, the market falls and rises, falls
and rises again, oscillating repeatedly as prices are gradually pulled up
fifty or sixty tawara above the floor. Then comes a brief spell of
indecision. After that, prices are pushed up relentlessly---ten tawara a
day---until yesterday's sellers are chasing to cover today, and what was
bought in the morning session returns in the afternoon at double the
volume. When people begin to clamor, that is the moment when prices
peak---the annual ceiling.

% --- Article 12 ---
\article{12}{On the Three Positions and the Kinoe Month: A Secret Tradition}

\begin{sokyubox}
{\small\itshape\color{sokyubar} \jp{第十二章\quad 三位甲月秘傳の事}}

\medskip
The ten Heavenly Stems (\jt{十干}{jikkan}) are divided as follows:

\textbf{\jp{甲}} (\textit{kinoe}), \textbf{\jp{乙}} (\textit{kinoto}),
\textbf{\jp{丙}} (\textit{hinoe}), \textbf{\jp{丁}}
(\textit{hinoto})---the first four stems. These S\=oky\=u associates
with rising markets.

\textbf{\jp{戊}} (\textit{tsuchinoe}), \textbf{\jp{己}}
(\textit{tsuchinoto}), \textbf{\jp{庚}} (\textit{kanoe}),
\textbf{\jp{辛}} (\textit{kanoto}), \textbf{\jp{壬}}
(\textit{mizunoe}), \textbf{\jp{癸}}
(\textit{mizunoto})---the remaining six stems. These he associates with
falling markets, with \textit{mizunoto} marking the floor.

Each of the twelve months of the calendar year is assigned a cyclical
stem. Depending on the year, the rise may begin in the \textit{kinoe}
month and conclude by the \textit{hinoto} month; or it may be delayed
into the \textit{tsuchinoe} or \textit{tsuchinoto} months. Eight or nine
times out of ten, the rise begins in the \textit{kinoe} month. When one
targets a hundred-tawara rise, there is no case in which one fails to
profit. In any event, within these four months, the market never fails
to rise. From \textit{tsuchinoe} through
\textit{tsuchinoto}---the
span of six months---there is no rise; these are months of decline, with
\textit{mizunoto} marking the month when the floor price emerges. There
are also cases, owing to unseasonable calamity or extraordinary events,
when the rise is delayed past \textit{tsuchinoe} and does not conclude
until the \textit{kanoto} month; but such an occurrence is exceedingly
rare---perhaps once in thirty years. The above has held true without
deviation from Meiwa~5 (1768) through Kansei~9 (1797).
\end{sokyubox}

\commentary
This section describes the method of predicting the monthly rise and
fall by means of the Heavenly Stems assigned to each month. The
practical application and historical examples are set out in later
chapters and at the end of the volume; the reader should consult those
passages.

Where the text says ``depending on the year, the rise begins in the
\textit{kinoe} month and concludes by the \textit{hinoto} month,'' and
so forth, it means that the rise commences in the \textit{kinoe} month
and is completed within the \textit{hinoto} month, after which prices
turn downward. The passages that follow should be understood on the same
principle.

\begin{sokyubox}
{\small\itshape\color{sokyubar} Article 12 (continued)}

\medskip
In a year when the seventh month falls on \textit{kinoe}, natural
disasters tend to strike both the \jt{上方}{Kamigata} region and this
locality, and prices frequently surge sharply from the seventh month
onward. In such a year, one should buy without hesitation from the sixth
month. Even if no disaster occurs, by the twelfth month the Osaka market
will have pulled prices up to a ceiling of a hundred tawara or more.
When the seventh month falls on \textit{kinoe}, one should buy regardless
of whether the harvest is good in both Kamigata and the local region.

Furthermore, in any month, if prices surge sharply by a hundred to a
hundred and twenty tawara in the \textit{kinoe} month, one should
understand that no further rise will follow---that is the ceiling. If the
rise is only fifty or sixty tawara, one should understand that further
upside remains in the \textit{hinoe} and \textit{hinoto} months.

In a year when the year-block (\jt{年塞}{toshifusagari}) turns toward
the west, a great rise is beyond doubt. Buy without hesitation from the
sixth month. The ceiling should appear in the eighth, ninth, tenth, or
eleventh month. Afterward, while prices consolidate thirty or forty
tawara below the ceiling, the \textit{kinoe} stem will rotate to the
third, first, or eleventh month---at which point another rise of fifty
or sixty tawara is beyond doubt. Such a year is one in which profits are
made on the autumn-to-spring oscillation.
\end{sokyubox}

\commentary
In any month to which \textit{kinoe} falls, if prices surge sharply by a
hundred to a hundred and twenty tawara in the \textit{kinoe} month,
there will be no further rise afterward---that becomes the ceiling and
prices turn downward. However, if the rise amounts to only fifty or
sixty tawara, that is insufficient for a ceiling, and prices should
continue rising through the \textit{hinoe} and \textit{hinoto}
months---two or three more months of advance.

Additionally, there is what is called the year-block
(\textit{toshifusagari}). The method of reckoning it is described at the
end of this volume. When the block turns toward the west, prices will
likewise rise without doubt through the eighth, ninth, tenth, and
eleventh months, so one should buy from the sixth month onward without
hesitation. After the ceiling appears around the eleventh month and
prices decline thirty or forty tawara into consolidation, the
\textit{kinoe} stem may happen to rotate to another month in the same
year or the following spring. At that time, another rise of fifty or
sixty tawara is assured; if one knows this, one can profit on the
autumn-to-spring oscillation. One should keep this in mind well in
advance. This section describes the historical experience of years when
the seventh month falls on \textit{kinoe}.

\begin{sokyubox}
{\small\itshape\color{sokyubar} Article 12 (continued)}

\medskip
When it comes to the ceiling price: even if in the first year the
ceiling reaches a hundred, a hundred and twenty, a hundred and sixty, or
even a hundred and eighty tawara, in the second year the ceiling will
fall short of the first year, and in the third year it will fall still
shorter. One should understand this. The decline the following spring
likewise falls short in proportion. The number of months and the number
of tawara, both rising and falling, should be studied in the price-scale
records (\jt{位付の書}{kuraizuke no sho}).

Even when the leading month of the Three Positions arrives and prices
fail to rise---indeed, even if they drop sharply by thirty or forty
tawara---one should not think this a mistake. A rise in the
\textit{hinoe} and \textit{hinoto} months is beyond doubt. When the
Three Positions rotate to the fifth, third, or first month, there is no
basis for estimating the new-crop price. Nevertheless, however good the
harvest that year, one must not sell short. If the price is in the
thirty-tawara range, one should buy.
\end{sokyubox}

\commentary
This section explains the general proportions when a ceiling price forms.

When the ceiling is struck, the typical pattern is as follows: even if in
the first year the ceiling rises by a hundred to a hundred and eighty
tawara, in the second year the proportional rise will fall short of the
first, and in the third year it will be still less. The spring decline
likewise diminishes year by year---the first year's decline exceeds the
second, and the second exceeds the third. The evidence for this is set
out plainly in the price records at the end of the volume; one should
study them and grasp the general pattern.

Also, even when the \textit{kinoe} month---the leading month of the
Three Positions---arrives and prices fail to rise, instead falling
sharply by thirty or forty tawara in the opposite direction, one should
not regard this as an error. A rise in the \textit{hinoe} and
\textit{hinoto} months is certain. However, when the Three-Position
cycle rotates to the fifth, third, or first month, forecasting is
difficult and speculation is hazardous. Even so, in such years one
should not take the sell side even if the harvest is good. If rice stands
at the thirty-tawara level, one should rather buy on speculation.

% --- Article 13 ---
\article{13}{On a Midwinter Ceiling Followed by a Sharp Drop in the Sixth Month}

\begin{sokyubox}
{\small\itshape\color{sokyubar} \jp{第十三章\quad 冬中天井六月急下げの事}}

\medskip
After a midwinter ceiling price, in a falling market: in a year of
four-and-a-half-tenths crop failure, market sentiment remains taut
through the following spring and prices do not decline. Then, as
shippers are in no hurry to buy, the midsummer \textit{doy\=o} period
brings clear skies, the paddies look splendid, shipments of domain rice
(\jt{御藏米}{okura-mai}) fall short, inbound shipping is scarce, and
sentiment turns sour---whereupon prices drop sharply in the sixth month.
\end{sokyubox}

\commentary
This chapter is of the same import as the passage in Chapter~9: it
describes a summer decline following a midwinter ceiling.

``Domain rice shipments fall short'' (\jp{御藏米出不足}) refers to a
decrease in exported rice, while ``shipping falls short'' (\jp{船不足})
means a shortage of inbound cargo---here it speaks of a scarcity of
buyers.

% --- Article 14 ---
\article{14}{On ``At Two, Close Out; at Three, Fully Ripe; at Four, It Turns''}

\begin{sokyubox}
{\small\itshape\color{sokyubar} \jp{第十四章\quad 二つ仕舞三つ十分四つ轉じの事}}

\medskip
After a ceiling price, the falling market typically declines by ten to
thirty or forty tawara each month over a span of five or six months.
However, depending on conditions, there are cases where the market
declines for three months and rises in the fourth; or where a further
decline begins in the fifth month; or where all six months are months of
decline.

The method is as follows: once the ceiling price has appeared, watch
carefully for five or six days to ten days. Even if prices hold at the
same level, or even if they are only slightly lower, sell. By the end of
that month, prices will certainly fall by twenty or thirty tawara. If,
from the first through the tenth of the following month, prices recover
five, six, or ten tawara above the prior month's low, then by that
month's end they will again fall twenty or thirty tawara. Pay close
attention to the first four or five days of each month. If volume is thin
and no trades materialize, continue selling through the tenth and into
the twentieth.

If, however, in the following month---from the first through the fourth
or fifth day, or even through the tenth---prices remain at the same level
as the prior month's low, or drop a further five, six, or ten tawara,
then that month's close will certainly reverse upward. One must understand
this thoroughly.

\begin{quote}
\textit{``At two, close out; at three, fully ripe; at four, it turns.''}\\
\jp{二つ仕舞、三つ十分、四つ轉じ}
\end{quote}

\noindent This is the secret tradition of the Three Positions.

The meaning is this: from the ceiling month, prices decline with
certainty for two or three months. In the fourth month, depending on
circumstances, prices may rise---it is dangerous, and one must not sell.
\end{sokyubox}

\commentary
This chapter explains the normal tendency of monthly decline after a
ceiling price, and simultaneously teaches the significance of the number
of months.

After the ceiling, the falling market declines each month by roughly
forty or fifty sen to a full yen over a span of five or six months.
However, depending on the occasion, prices may fall continuously for
three months and then rise from the fourth month onward; or a further
decline may begin in the fifth month; or all six months may be months of
decline.

The method of selling in such a market is this: after the ceiling price
appears, watch carefully for five or six to ten days, then sell at the
second peak---that is, at the ``second mountain'' (\textit{ni no yama}).
If no second mountain forms and the market merely consolidates, sell even
if prices appear somewhat low. By month's end, the decline should amount
to forty or fifty sen to a full yen. Then, owing to seller profit-taking
(\jt{利喰}{rigui}) and the buying enthusiasm of the new session opening
(\jt{發會}{hakkai}), if prices recover somewhat through the tenth of the
following month, they should again decline by forty or fifty sen to a
full yen by that month's end.

In short, the pattern of a falling market is: the month opens high and
closes low. One should focus on selling in the first four or five days.
If volume is thin and one cannot sell sufficiently, one should continue
selling through the tenth or even the twentieth. This pattern of
oscillating downward---rallying then declining, rallying then
declining---is the normal tendency of a falling market.

However, if this normal pattern breaks down---if prices consolidate at
the same level as the previous month's low, or if at the month's opening
they drop sharply by forty or fifty sen to a full yen---one must take
careful note. This may mark the floor, and by month's end prices will
certainly rise. This is described with reference to the pattern of three
months of decline followed by a rise in the fourth month. (Note: the
prices used in this commentary are set for the convenience of explanation
and should not be taken literally. The same applies to many examples
hereafter.)

\begin{sokyubox}
{\small\itshape\color{sokyubar} Article 14 (continued)}

\medskip
In any month, if from the previous month's low, prices are ten tawara
higher by the first through the twentieth of the following month, they
will decline twenty tawara by that month's end. If, conversely, prices
fall five, six, or ten tawara below the previous month's low at the
start of the following month, prices will certainly rise by the latter
part of that month. This is the art of profiting where others do not
see. Consider it well and put it to use; study the price-scale records
and learn from them.
\end{sokyubox}

\begin{sokyubox}
{\small\itshape\color{sokyubar} Article 14 (continued)}

\medskip
The above works as follows: sell at the opening of the month; when
prices fall at month's end, buy back the sell position and, in addition,
go long. Wait for the rise at the opening of the following month, then
sell again that month, and also close the long position. In this manner,
trade briskly for two or three months. To carry a sell position for five
or six months at a stretch is insufficient in profit and exceedingly
dangerous---one must not sell on extended terms. This period is the
selling opportunity of the entire year. Apart from this, to sell at other
times invariably results in loss---this is plain, and one must exercise
caution. The above pertains to rice in a year of great crop failure, with
prices up more than a hundred tawara---a year of a major rise.
\end{sokyubox}

\commentary
This section continues the preceding passage, describing the normal
pattern of a falling market and the trader's method.

In any event, when the following month arrives and prices recover by
forty or fifty sen above the previous month's floor by the fifteenth or
twentieth, they should fall by roughly double---a full yen---by month's
end. Conversely, when prices decline by thirty or fifty sen at the
month's opening, they will certainly rise by month's end. Because few
recognize this rhythm, one can profit as follows: sell at the month's
opening, reverse into a long position (\textit{doten}) at the month-end
decline, then reverse again into a sell at the following month's opening
rally---trading in this manner, one can profit skillfully where others do
not see.

However, one important caution: this method of trading is limited to two
or three months. To continue it for five or six months is forbidden.

In sum, the selling opportunity of the entire year is confined to the
period after the annual ceiling price, as described above. Those who seek
certain victory must not engage in sell-side speculation outside of this
window---one must exercise the greatest caution.

% --- Article 15 ---
\article{15}{Rise and Fall Are Matters of Natural Order}

\begin{sokyubox}
{\small\itshape\color{sokyubar} \jp{第十五章\quad 高下は天性自然の事}}

\medskip
The rise and fall of rice prices is governed by the principle of natural
order (\jt{天性自然}{tensei shizen}). Because prices move according to
this natural principle, it is exceedingly difficult to determine with
certainty that they will rise or fall. Those unfamiliar with this Way
must not engage carelessly in this trade.
\end{sokyubox}

\commentary
The rise and fall of the rice market is, in the final analysis, governed
by natural order---cheap in bountiful years, dear in years of famine.
Yet the timing of these movements---that is, the moment of change---is
influenced now by the conditions of rice supply and distribution, now by
favorable or adverse weather, now by speculative reasoning about the
causes of abundance or scarcity. The causes that move the market shift
from moment to moment; and moreover, there are times when the market
responds to these causes slowly and times when it responds quickly. For
this reason, one cannot pronounce with certainty that prices will
necessarily rise or necessarily fall.

\jt{慈雲齋}{Jiunsai} too has said: \textit{``Between reason and unreason
lies a reason beyond reason---know this as the wellspring of rise and
fall.''} And yet, in the end, there can be no such thing as a ``reason
beyond reason'' in this world. When our intellect is insufficient to
grasp the principle, we provisionally call it ``reason beyond reason'' or
``the realm of the unfathomable.'' So it is with the rise and fall of
markets: each rise has its cause, and each fall has its cause. To invoke
``reason beyond reason'' is ultimately to confess that one's investigation
of causes has not yet reached its conclusion. Yet to pursue the full
chain of cause and response is supremely difficult, and the difficulty of
predicting market movements stems, in the end, from the inexhaustible
variety of these causal states.
